% !PRINTED 450
\weberpage{450}
%<<CONT
may be entirely independent of one another; but the case is not excluded that certain relations are fixed between them, which are to be taken into account in calculation.

A function of a function field counts as zero only when it vanishes identically, that is, for all values of the variables that come into consideration.

Since for the present we do not wish to restrict our considerations, we shall now speak of \gesperrt{fields} simply, and shall designate the objects with which the calculations are to be carried out, which may be numbers as well as functions, as \gesperrt{quantities} or also as the \gesperrt{elements} of the field $^1$).

A \gesperrt{field} is then a system of quantities of such completeness that in it the quantities can be added, subtracted, multiplied and divided.

When all the quantities of one field are contained in a second field, the first field is called a \gesperrt{divisor} of the second.

If $\Omega$ is a divisor of $\Omega'$, we shall also say that $\Omega'$ \gesperrt{is a field over} $\Omega$.

% [sic] The German source prints the anomalous hyphenated phrase "ein- für allemal"; the idiomatic translation is retained.
Strictly speaking, the single number ``zero'' forms a field by itself. But for the sake of brevity we shall exclude this once and for all.

Then the most obvious example of a field is the \gesperrt{totality of all rational numbers}.

This field is a divisor of every other field; for every field contains at least one quantity $\omega$ different from zero, hence also the quotient $\omega:\omega=1$, that is, the number 1; hence also, since all integers arise by addition and subtraction of units, all integers; and from the integers one can again derive all fractions by division.

Other examples of number fields are: the totality of all complex numbers of the form $x+yi$, in which $i=\sqrt{-1}$, and $x$ and $y$ denote all rational numbers; further the totality of all (rational and irrational) \gesperrt{real numbers}, or the totality of all complex numbers $x+yi$ that exist at all.

As an example of a function field may serve the totality of all integral and fractional rational functions of a
%CONT>>

\par\noindent\rule{4cm}{0.4pt}\par
{\footnotesize $^1$) Cf. the author's paper: ``Die allgemeinen Grundlagen der \gesperrt{Galois}'schen Gleichungstheorie.'' Mathematische Annalen, Bd.\,43.}
